\section{Evaluation framework}
\label{sec:evaluation}

Roughly half of what this project has produced is its evaluation method, and it is presented here as
method rather than as process, because each component was adopted in response to a specific measured
failure and each is portable to another game. \S\ref{sec:background-vsix} describes the failures.

\subsection{The outcome measure}
\label{sec:evaluation-measure}

The primary outcome is the team win rate, the proportion of games in which the evaluated
configuration's team takes at least five of the nine half-suits. Comparisons are reported as an
\textbf{edge in percentage points},
\[
  \text{edge} \;=\; 100 \times (\text{win rate} - 0.5),
\]
so that zero denotes indistinguishable configurations and the sign gives the direction. The paper
writes \emph{pp} for percentage points of this quantity. Three other quantities appear and are never
conflated with it: the mean number of half-suits a team takes, used for the search's internal return;
predictive log loss and argmax hit rate, used to score posteriors as predictors
(\S\ref{sec:background-vsix}); and the adversarial edge of \S\ref{sec:results-adversarial}, which is
an edge in the same units but measured for a fitted opponent against the evaluated configuration.

\subsection{Duplicate deal blocks}
\label{sec:evaluation-duplicate}

A \emph{cell} compares two configurations $A$ and $B$ over a \emph{bank} of deals identified by one
integer seed. Every deal in the bank is played twice: once with $A$ in the odd seats and $B$ in the
even seats, and once with the seats exchanged. Both configurations therefore play the identical deal
from both sides.

Fish has large per-deal variance, since a deal can be won or lost on the cards, and an unpaired
comparison spends most of its sample resolving that variance rather than the difference between
policies. The rotation cancels the deal's intrinsic advantage exactly, in the spirit of duplicate
bridge and of the variance-reduction estimators developed for poker
\cite{zinkevich-divat,burch-aivat}. The design requires reproducibility from the seed alone
(\S\ref{sec:inference-engine}); without it the two rotations are not the same deal and the design
degrades silently into an unpaired one.

\subsection{Deal-clustered bootstrap confidence intervals}
\label{sec:evaluation-intervals}

The two rotations of one deal are not independent observations. They share a deal and their outcomes
are negatively correlated by construction, so treating them as independent understates uncertainty.

Every interval in this paper is a two-sided 95\% \textbf{deal-clustered bootstrap confidence
interval}. The resampling unit is the \emph{deal}, carrying both of its rotations together; 20{,}000
bootstrap resamples are drawn with replacement from the deals of the cell, the edge is recomputed on
each resample, and the interval is the empirical 2.5th and 97.5th percentiles of that distribution.
No Wilson or normal-approximation interval on the game count is quoted anywhere in this paper, and
the substitution is not cosmetic: at these sample sizes it is the difference between a claim that
replicates and one that does not.

Two derived quantities are used throughout and are defined once here.

\begin{description}
\item[Pooling two banks.] Two banks of equal size pool as the mean of their two edges, with standard
  error $\mathrm{se}_{\text{pooled}} = \sqrt{\mathrm{se}_1^2 + \mathrm{se}_2^2}/2$, where each
  $\mathrm{se}$ is recovered from that bank's clustered interval as half-width divided by $1.96$. The
  pooled interval is the estimate plus or minus $1.96\,\mathrm{se}_{\text{pooled}}$. Both per-bank
  values are printed beside every pooled figure in this paper.
\item[Differences between separately measured cells.] A difference between two pooled cells, such as
  an ablation's leave-one-out drop, is the difference of the estimates with the two pooled half-widths
  combined in quadrature. This is conservative. The configurations are paired on deals, so a genuine
  paired difference would be narrower, but the harness produces no paired difference \emph{across} two
  separately executed cells, and constructing one would assert more than the data structure supports.
\end{description}

\subsection{Power computed before a cell is run}
\label{sec:evaluation-power}

At parity a cell's 95\% half-width in percentage points is approximately $98/\sqrt{N}$ for $N$ games.
The harness computes this for every cell and emits it with the result. The arithmetic is applied in
advance because it determines what may be attempted: resolving one point requires roughly
\vsevenScoreboardGames{} games, and a quarter of a point requires sixteen times that.

The failure this prevents is the one \S\ref{sec:background-vsix} describes, in which five mechanisms
cleared a 95\% paired interval at one budget and returned exactly zero at three times it. Sample
sizes in this paper are chosen from the arithmetic before a cell is run, never after seeing its
result.

\subsection{Replication across disjoint banks as an advance-specified criterion}
\label{sec:evaluation-replication}

Every claim is measured on two disjoint deal banks, and the criterion was fixed before the data were
seen:

\begin{quote}
A claim whose \textbf{sign} does not agree on both banks is reported as not replicated, whatever its
pooled interval says.
\end{quote}

This is stricter than a pooled significance test, deliberately. A pooled interval can exclude zero
while the two banks disagree, and that configuration has occurred in this corpus and has previously
been published as a positive result. The criterion binds against this paper's own findings in
\S\ref{sec:results-attribution} and again in \S\ref{sec:results-panel}, and both instances are
reported.

\subsection{Sealed material and advance registration}
\label{sec:evaluation-seal}

A programme that generates candidate configurations, scores them and selects the best is performing a
maximisation, and the selected configuration's score is upward biased by the amount a maximum over
$K$ draws exceeds the mean. No care taken in the final measurement removes that bias if the final
measurement uses material the selection has seen.

Before any candidate existed, deal banks were reserved by seed and their digests committed, a digest
being a rolling hash over the six dealt hands and the dealer of every deal in the bank, computed
without constructing a policy or playing a game. An adversary population was split in half by a rule
fixed in advance, and one half was encrypted with the plaintext's SHA-256 committed, so that it could
not be trained against even accidentally. The engine enforces the seal: a run naming a sealed seed
fails unless an explicit environment variable is set, which was set once, at the start of the final
evaluation, and recorded.

The entire final battery, comprising every cell, its sample size, its threshold, the arithmetic
behind the sample size and the conditions that would count as non-confirmation, was written down and
committed before any holdout bank had been played and before the sealed half had been decoded. The
evaluation was then executed from a cleared context, reading that document and nothing else from the
project's records. Where it required a training figure for comparison it took it from the registered
protocol, which states such expectations in advance for exactly this purpose, so that agreement
constitutes a check and disagreement constitutes a finding.

Anything the evaluation did that the protocol does not specify is recorded as a deviation, next to
the affected result. There are \vsevenDeviations{}, listed in Appendix~\ref{app:deviations}.

\subsection{Selection effects and multiple comparisons}
\label{sec:evaluation-multiplicity}

This study evaluates many candidate mechanisms, opponents, ablations, adversarial searches and
configurations, and the inferential status of that needs stating precisely.

\textbf{No formal family-wise or false-discovery correction was applied}, and no interval reported
here is adjusted for multiplicity. What the design provides instead is structural.
\emph{Registration} fixes the primary comparison, its threshold and its interpretation before the
data exist, so the primary claim is a single pre-specified test rather than the best of many.
\emph{Sealing} makes the holdout material unavailable to the selection process by construction rather
than by discipline. \emph{Disjoint banks with a sign criterion} require a claim to appear twice
before it is reported. \emph{Calibration} of the detection floor (\S\ref{sec:evaluation-floor})
establishes what the instrument can resolve, so that a null is interpretable. And an explicit
\emph{selection check} (\S\ref{sec:results-selection}) computes how much of the measured gain a
maximum over $K$ candidates would produce under the null and compares advance-stated training figures
against their holdout counterparts.

Exploratory comparisons in \S\ref{sec:development} are reported as exploratory: they were run on
training material during the development phase, none is treated as confirmatory, and none is used to
support the paper's primary claim.

\subsection{A soundness gate applied before any strength measurement}
\label{sec:evaluation-gate}

Win rate and soundness are dissociable in this policy family. The corpus contains a configuration
that scores \vsevenScoreboardTrapWin\% against an earlier version while making
\vsevenScoreboardTrapDead\% provably dead asks, asks for cards its own team demonstrably already
owns, and losing \vsevenScoreboardTrapLimit\% of its games to the action cap. It is both stronger and
broken.

A configuration is therefore gated \emph{before} any strength number is computed for it, on \vsevenGateRulesWord{}
rules covering dead asks, dead runs, action-limit games, the mirror game-length tail, late
declarations, and the side-channel tests below. A configuration that fails is rejected regardless of
its win rate. The gate runs on \emph{training} material by design, being a soundness check on
self-play behaviour rather than a strength claim. It is accompanied by a negative control, a
configuration constructed to be rejected, because a gate that cannot reject it would certify nothing
about the configuration it accepts.

\subsection{Calibrated detection floors}
\label{sec:evaluation-floor}

An exploitability figure is produced by a search, and a search that finds nothing admits two
explanations: there is nothing to find, or the search is too weak. A \textbf{detection floor}
separates them. An artificial handicap of known size is planted in a target, a responder is fitted
against the handicapped target, and the \emph{excess} of that responder's edge over its edge against
the unhandicapped target is measured. The smallest planted edge the procedure recovers is the floor.

Two findings from that calibration shape every exploitability statement in this paper. First, the
floor does not decrease with evaluation games. An early prediction that it would fall as
(evaluation games)$^{-1/2}$ was refuted at its own predicted value: at four times the games the floor
is \vsevenPrimaryFloor{}~pp, and the sub-floor rung remains undetected not because the interval is
wide but because the measured excess resolves \emph{to} zero. Below roughly a point and a half a
fitted responder is empty rather than merely unresolved, so the binding constraint is search power
rather than evaluation power. Second, floors are family-specific: the declaration channel has its own
and higher floor of \vsevenDeclFloor{}~pp, and reading a declaration-objective result against the
general floor would overstate it.

An adversary class contributes to an exploitability statement only after passing a planted-edge
calibration at or below the effect size in question. That precondition is why the negative controls of
\S\ref{sec:results-controls} are run: an instrument that cannot recover a planted edge of a given size
cannot report the absence of a real one of that size.

\subsection{Mechanical homogeneity and side-channel controls}
\label{sec:evaluation-side}

The evaluated team is three identical agents with no shared secret. That their coordination is common
knowledge is legitimate (\S\ref{sec:game-why}), but it also means an implementation detail could
become an unintended signalling channel. Before this cycle the constraint had only ever been checked
by reading the source.

Four mechanical tests now run over all four engine decision types. \textbf{S3, listening
substitution} asks whether a seat's behaviour changes more when a teammate's hidden action changes
than when an opponent's does, beyond what the rules explain. \textbf{S4, stream independence} asks
whether the seat's private random stream is drawn independently of the deal. \textbf{S5, posterior
invariance} asks whether the seat's action depends on true card locations beyond what its own
posterior supports. \textbf{S6, seat isolation} replays every decision from own hand, public event
stream, rules and reset seed alone, and counts the decisions the reconstruction cannot reproduce, at
zero tolerance.

The tests are calibrated by three positive controls, each a policy carrying one planted violation,
and each must fail exactly the tests the protocol names for it in advance and no others.
\S\ref{sec:design-material} reports the calibration on holdout material. Without those controls a
clean report from the gate would be uninformative.

One limitation is stated wherever the gate is used. The zero-tolerance form of S6 requires a single
worker thread, because the audit depends on the deal-to-thread partition and work stealing does not
fix that partition at higher counts, and it requires agents rebuilt for each deal. Rebuilding agents
per deal \emph{changes play}. The gate therefore certifies that with cross-deal residue removed,
nothing else depends on anything it should not. The residue itself is not removed from the deployed
mode, and every strength number in this paper is measured under it
(\S\ref{sec:results-controls}).

\subsection{Worst case and minimax regret over a shared panel}
\label{sec:evaluation-panel}

The project's standing reporting rule is that the worst case leads and an aggregate never does, since
a mean over an opponent panel rewards beating weak opponents by larger margins, which is not the
quantity of interest. Two statistics are reported over a panel scored identically for every
configuration: the \textbf{worst cell}, being the panel member on which a configuration does least
well, and the \textbf{minimax regret}, being for each member the shortfall from the best configuration
on that member, maximised over members. They can disagree, and in \S\ref{sec:results-panel} they do.
Reporting both, and reporting the panel mean last and labelled as a diagnostic, is the discipline.

Every configuration must carry a cell for every member or minimax regret is undefined, which is why
the panel is fixed by name in advance and scored for all configurations on the same banks. One
consequence of that design is recorded where it matters: a configuration that is itself a panel
member meets itself and carries a free zero cell, which can only flatter its worst case, so
\S\ref{sec:results-panel} prints the worst case both with and without self-cells.
